\section{Исходные выпрямления и гармонические семейства}
\subsection{Обозначения и проверенная исходная конструкция}
Поле $k$ имеет характеристику нуль. Полиномы Чебышёва нормированы как
\[
C_j(u)=2T_j(u/2),\qquad C_0=2,\quad C_1=u,\quad C_{j+1}=uC_j-C_{j-1}.
\]
Рассматриваются замкнутые вложения, то есть $k[f_1,f_2,f_3]=k[u]$.
Полином параметра $P$ удовлетворяет $P\circ\gamma=u$; координата, кроме
того, должна дополняться до полиномиального автоморфизма. Эти два
условия не равнозначны.

Исходный препринт \cite{Repository} содержит автоморфизмы
$\Phi_T,\Phi_H$ и их полиномиальные обратные над
$K=\Q[a]/(a^2-a+1)\simeq\Q(\sqrt{-3})$ для
\[
\begin{aligned}
\gamma_T(u)&=(u^3-3u,u^4-4u^2,u^5-10u),\\
\gamma_H(u)&=(C_3(u),C_5(u),C_7(u)).
\end{aligned}
\]
Проверены $\Phi_T\gamma_T=\Phi_H\gamma_H=(u,0,0)$, обе обратные
композиции посредством точных полиномиальных тождеств и все двенадцать
напечатанных компонент. Степени компонент равны соответственно
$(5,15,24)$ и $(7,21,41)$; их минимальность не утверждается.
Полные формулы сохранены в исходном препринте и в автономном
вычислительном пакете. Замена
\[
M(x,y,z)=(x,y+2,z-5x)
\]
переводит $\gamma_T$ в $(C_3,C_4,C_5)$.

\subsection{Конструктивный перенос через модификацию}
Пусть $A=(x,H,W)$ и $A^{-1}=(x,\mathcal Y,\mathcal Z)$ ---
взаимно обратные автоморфизмы, $m\in k[x]$, $\deg m>0$, $p\in k[x,y]$.
Рассмотрим $\beta=(x,p(x,y)+m(x)z,y)$. В факторкольце
$S=k[x]/(m)$ положим $c(y)=H(x,p(x,y),y)$.
Если линейный коэффициент $c$ обратим, а все старшие коэффициенты
нильпотентны, то подстановка $y\mapsto c(y)$ задаёт автоморфизм $S[y]$.
Его композиционный обратный $\psi$ строится
конечным подъёмом по нильпотентному идеалу старших коэффициентов.

Выбираем полиномиальный представитель
\[
\varphi(x,T)\equiv W(x,p(x,\psi(T)),\psi(T))\pmod m.
\]
Тогда
\[
\mathcal A=\left(x,H\circ\beta,
\frac{W\circ\beta-\varphi(x,H\circ\beta)}m\right)
\]
и
\[
\mathcal A^{-1}(x,T,Q)=
\left(x,\mathcal Z(x,T,mQ+\varphi),
\frac{\mathcal Y(x,T,mQ+\varphi)-p(x,\mathcal Z(x,T,mQ+\varphi))}m\right)
\]
имеют полиномиальные компоненты: оба числителя делятся на полный $m$,
включая кратности его корней. После обращения $m$ обе композиции
тождественны. Инъективность локализации переносит равенства в исходное
кольцо, включая исключительный слой $m=0$.
Это конкретная реализация классических механизмов подъёма и относительных
координат; общая идея переноса не объявляется новой \cite{EVM,Repository}.
Случай $m\in k^*$ проще: сама $\beta$ уже является автоморфизмом.

Если $H\circ\gamma=u$ и $w=W\circ\gamma$, то из относительной пары
получается полный выпрямитель
\[
\Phi=(H,\ x-f_1(H),\ W-w(H)).
\]
Наличие только полинома параметра, без проверенной координатной пары,
для этого вывода недостаточно.

\subsection{Полиномиальные изменения гармонических индексов}
Для взаимно простых $a,b$ и знаков $\varepsilon_a,\varepsilon_b\in\{\pm1\}$
условия
\[
c'\equiv\varepsilon_a c\pmod{2a},\qquad
c'\equiv\varepsilon_b c\pmod{2b}
\]
дают
\[
C_{c'}-\varepsilon_a\varepsilon_b C_c\in\Q[C_a,C_b].
\]
Следовательно, соответствующие пространственные кривые связаны
полиномиальным сдвигом третьей координаты.

Краткое доказательство использует узлы плоской проекции
$(C_a,C_b)$. Её кольцо состоит ровно из полиномов по $u$, принимающих
одинаковые значения на двух прообразах каждого узла: включение очевидно,
а равенство следует сравнением коразмерностей
$(a-1)(b-1)/2$ через полугруппу $\langle a,b\rangle$ и интерполяцию.
Указанные сравнения индексов дают требуемые равенства значений по
тригонометрической формуле Чебышёва. Исключение ведущих членов даёт
рациональный полином сдвига. Подробности и точные примеры содержатся
в исходном аппендиксе; литературная основа узловой проекции --- \cite{FF}.

Отсюда следуют два полных семейства над $K$:
\[
H(3,4,c),\quad\gcd(c,12)=1;
\qquad H(3,5,c),\quad\gcd(c,15)=1.
\]
Для пары $(3,4)$ достаточно представителей $c=1,5$ по модулю $24$;
для $(3,5)$ --- $c=1,2,4,7$ по модулю $30$. Представители с индексом
$1$ очевидны, с индексом $2$ известны классически, $c=4$ сводится
перестановкой к $(3,4,5)$, а оставшиеся два обеспечены исходными
выпрямлениями.

Индекс $1$, допустимые модели $H(2,b,c)$ и $H(a,b,a+b)$ уже были
известны в литературе \cite{FF}; они не приписываются этой работе.
Топологическая классификация гармонических узлов \cite{KP} также не
является классификацией комплексных вложений по автоморфизмам.

\subsection{Конечный перебор и его точная область}
Перебраны $44142$ попарно взаимно простые тройки
$2\le a<b<c\le100$ по всем строго убывающим элементарным редукциям
\[
c=ra+sb\quad\longmapsto\quad |ra-sb|,
\qquad (x,y,z)\longmapsto(x,y,C_r(x)C_s(y)-z).
\]
Из них $3209$ достигают индекса $1$, ещё $237$ достигают проверенной
базовой модели, а $40696$ остаются неразрешёнными \emph{этим алгоритмом}.
Последнее число не является числом невыпрямляемых кривых.

Далее для любого вложения $\gamma=(f_1,\ldots,f_n)$ используется
\[
A=k[t,u,X_1,\ldots,X_n],\qquad
D_\mu=t^\mu\partial_u+\sum_i f_i'(u)\partial_{X_i},\qquad
V=\gamma(u)-t^\mu X,
\]
где $\mu\ge1$. Это LND на $n+2$ переменных. Его нельзя смешивать
с якобианным LND модификации на трёх переменных.
